% GENERATED FILE. Edit canonical lessons, then run npm run generate:books.
% canonical-source-hash: fnv1a64:16f3c888ab5e9751
% canonical-lessons: ZH-C290-xinwen, ZH-C290-baodao, ZH-C290-diantai, ZH-C290-zhuchiren, ZH-C290-redian

\chapter{In the News}
\label{ch:zh-in-the-news}
\hlchaptermodality{290}

\begin{chapteropening}
\textbf{By the end of this chapter:} \emph{I can talk about the news, the radio, the host of a show and what everyone is talking about.}

Complete the last lesson of chapter 290: I can talk about the news, the radio, the host of a show and what everyone is talking about.
\end{chapteropening}

\section[xīnwén]{xīnwén — the news; news on TV and radio}

\label{lesson:ZH-C290-xinwen}



\begin{quote}
Before the new one: say the Mandarin for if, then the Mandarin for otherwise.
\end{quote}

\subsection*{You'll want to know: xīnwén}
\emph{xīnwén} — \textquotedblleft{}the news; news on TV and radio\textquotedblright{}.

\begin{quote}
\emph{wǒ měitiān tīng xīnwén} — I listen to the news every day.
\end{quote}

Its characters come later; for now it is a word for the ear.

\begin{cousinweb}
\emph{xīnwén} is built from \emph{xīn} 'new' and \emph{wén} 'to hear', both inherited from Middle Chinese, \emph{sin} and \emph{mjun} (Guangyun, in Baxter's transcription), which glosses the first 'new, as against old' and the second 'to know by its sound': what is newly heard.
\end{cousinweb}

\begin{grammarlens}[title={What you've built}]
One more word for telling what happened.
\end{grammarlens}

\subsection*{Your turn}
\begin{itemize}
\raggedright
  \item \emph{Say it:} \emph{xīnwén}
  \item \emph{Say it:} \emph{xīnwén}, once more
  \item \emph{Say it:} \emph{wǒ měitiān tīng xīnwén} — I listen to the news every day.
  \item \emph{Recall:} say the Mandarin for if, then the Mandarin for otherwise, then say \emph{xīnwén} again
\end{itemize}

\begin{tcolorbox}[breakable,colback=teal!4,colframe=teal!35!black,title={Before you move on}]
What does xīnwén mean? (\textquotedblleft{}The news; news on TV and radio\textquotedblright{}.) Say it once more.
\end{tcolorbox}
\section[bàodào]{bàodào — to report on; a news report}

\label{lesson:ZH-C290-baodao}



\begin{quote}
Before the new one: say the Mandarin for otherwise, then the Mandarin for the news on TV.
\end{quote}

\subsection*{You'll want to know: bàodào}
\emph{bàodào} — \textquotedblleft{}to report on; a news report\textquotedblright{}.

\begin{quote}
\emph{diànshì bàodào le zhège xiāoxi} — The news was reported on TV.
\end{quote}

Its characters come later; for now it is a word for the ear.

\begin{cousinweb}
\emph{bàodào} is built from \emph{bào} 'to report' and \emph{dào}, both inherited from Middle Chinese, \emph{pawH} and \emph{dawX} (Guangyun, in Baxter's transcription), which glosses the first 'to report' and the second 'a road' and 'a principle'; its sense 'to say' is not given there.
\end{cousinweb}

\begin{grammarlens}[title={What you've built}]
One more word for telling what happened.
\end{grammarlens}

\subsection*{Your turn}
\begin{itemize}
\raggedright
  \item \emph{Say it:} \emph{bàodào}
  \item \emph{Say it:} \emph{bàodào}, once more
  \item \emph{Say it:} \emph{diànshì bàodào le zhège xiāoxi} — The news was reported on TV.
  \item \emph{Recall:} say the Mandarin for otherwise, then the Mandarin for the news on TV, then say \emph{bàodào} again
\end{itemize}

\begin{tcolorbox}[breakable,colback=teal!4,colframe=teal!35!black,title={Before you move on}]
What does bàodào mean? (\textquotedblleft{}To report on; a news report\textquotedblright{}.) Say it once more.
\end{tcolorbox}
\section[diàntái]{diàntái — radio station}

\label{lesson:ZH-C290-diantai}



\begin{quote}
Before the new one: say the Mandarin for the news on TV, then the Mandarin for to report on the news.
\end{quote}

\subsection*{You'll want to know: diàntái}
\emph{diàntái} — \textquotedblleft{}radio station\textquotedblright{}.

\begin{quote}
\emph{wǒ bàba zài diàntái gōngzuò guo} — My dad once worked at a radio station.
\end{quote}

Its characters come later; for now it is a word for the ear.

\begin{cousinweb}
\emph{diàntái} is built from \zh{电} 'electricity' and \emph{tái} 'a platform', both inherited from Middle Chinese, \emph{denH} and \emph{doj} (Guangyun, in Baxter's transcription), which glosses the first as lightning and the second 'a high, square mound of earth'.
\end{cousinweb}

\begin{grammarlens}[title={What you've built}]
One more word for telling what happened.
\end{grammarlens}

\subsection*{Your turn}
\begin{itemize}
\raggedright
  \item \emph{Say it:} \emph{diàntái}
  \item \emph{Say it:} \emph{diàntái}, once more
  \item \emph{Say it:} \emph{wǒ bàba zài diàntái gōngzuò guo} — My dad once worked at a radio station.
  \item \emph{Recall:} say the Mandarin for the news on TV, then the Mandarin for to report on the news, then say \emph{diàntái} again
\end{itemize}

\begin{tcolorbox}[breakable,colback=teal!4,colframe=teal!35!black,title={Before you move on}]
What does diàntái mean? (\textquotedblleft{}Radio station\textquotedblright{}.) Say it once more.
\end{tcolorbox}
\section[zhǔchírén]{zhǔchírén — presenter; host of a show}

\label{lesson:ZH-C290-zhuchiren}



\begin{quote}
Before the new one: say the Mandarin for to report on the news, then the Mandarin for a radio station.
\end{quote}

\subsection*{You'll want to know: zhǔchírén}
\emph{zhǔchírén} — \textquotedblleft{}presenter; host of a show\textquotedblright{}.

\begin{quote}
\emph{zhège jiémù de zhǔchírén hěn yǒuyìsi} — The host of this show is a lot of fun.
\end{quote}

Its characters come later; for now it is a word for the ear.

\begin{cousinweb}
\emph{zhǔchírén} is built from \emph{zhǔ} 'to be in charge', \emph{chí} 'to hold' and \zh{人} 'a person', all inherited from Middle Chinese, \emph{tsyuX}, \emph{dri} and \emph{nyin} (Guangyun, in Baxter's transcription), which glosses the first 'to be in charge of' and the second 'to hold': the person who holds the show together.
\end{cousinweb}

\begin{grammarlens}[title={What you've built}]
One more word for telling what happened.
\end{grammarlens}

\subsection*{Your turn}
\begin{itemize}
\raggedright
  \item \emph{Say it:} \emph{zhǔchírén}
  \item \emph{Say it:} \emph{zhǔchírén}, once more
  \item \emph{Say it:} \emph{zhège jiémù de zhǔchírén hěn yǒuyìsi} — The host of this show is a lot of fun.
  \item \emph{Recall:} say the Mandarin for to report on the news, then the Mandarin for a radio station, then say \emph{zhǔchírén} again
\end{itemize}

\begin{tcolorbox}[breakable,colback=teal!4,colframe=teal!35!black,title={Before you move on}]
What does zhǔchírén mean? (\textquotedblleft{}Presenter; host of a show\textquotedblright{}.) Say it once more.
\end{tcolorbox}
\section[\texorpdfstring{\zh{热点} (rèdiǎn)}{热点 (rèdiǎn)}]{\zh{热点} (rèdiǎn) — hot topic; what everyone is talking about}

\label{lesson:ZH-C290-redian}



\begin{quote}
Before the new one: say the Mandarin for a radio station, then the Mandarin for the presenter of a show.
\end{quote}

\subsection*{You'll want to know: \zh{热点}}
\textbf{\zh{热点}} — \emph{rèdiǎn} — \textquotedblleft{}hot topic; what everyone is talking about\textquotedblright{}.

\begin{quote}
\textbf{\zh{这是今天的热点}} \emph{zhè shì jīntiān de rèdiǎn} — This is today's hot topic.
\end{quote}

\begin{cousinweb}
\textbf{\zh{热点}} (\emph{rèdiǎn}) is built from \zh{热} 'hot' and \zh{点} 'a dot', both inherited from Middle Chinese, \emph{nyet} and \emph{temX} (Guangyun, in Baxter's transcription), which glosses the first 'burning, as by fire' and the second 'a dot or stroke': a hot spot, the story everyone is talking about.
\end{cousinweb}

\begin{grammarlens}[title={What you've built}]
One more word for telling what happened.
\end{grammarlens}

\subsection*{Your turn}
\begin{itemize}
\raggedright
  \item \emph{Say it:} \emph{rèdiǎn}
  \item \emph{Say it:} \emph{rèdiǎn}, once more
  \item \emph{Say it:} \emph{zhè shì jīntiān de rèdiǎn} — This is today's hot topic.
  \item \emph{Recall:} say the Mandarin for a radio station, then the Mandarin for the presenter of a show, then say \emph{rèdiǎn} again
\end{itemize}

\begin{tcolorbox}[breakable,colback=teal!4,colframe=teal!35!black,title={Before you move on}]
What does \zh{热点} mean? (\textquotedblleft{}Hot topic; what everyone is talking about\textquotedblright{}.) Say it once more.
\end{tcolorbox}
